% Complete PS proof; independently reviewed by both continuation peers.
\paragraph{C.3bw. Constant proper parameter surfaces (Proved).}
Let $k$ be algebraically closed, let $(R,\mathfrak m)$ be a
noetherian local coefficient-$k$ algebra with residue field $k$,
and let $S/k$ be a smooth connected projective surface. We use
perfect $K$-theory and put $B=\operatorname{Spec}R\times_k S$.
Then, integrally,
\begin{equation}
 F_\gamma^3K_0(B)=0,\qquad
 A(\psi^2)=0\text{ on }K_0(B),\qquad
 A(X)=(X-1)(X-2)(X-4).
                                                        \tag{PS.1}
\end{equation}
These assertions do not require regularity, normality,
henselianity, or a characteristic hypothesis on $R$.

Suppose additionally that $\operatorname{char}k=p>0$ and that
$g:Z\to\operatorname{Spec}R$ is a proper morphism from a regular
noetherian scheme, an isomorphism away from the closed point.
Put $E=(Z\times_Rk)_{\mathrm{red}}$. For every
$\eta\in K_1(E\times_kS)_{\mathbf Q}$, the entire class
\begin{equation}
                    A(\psi^2)\eta
                                                        \tag{PS.2}
\end{equation}
lifts to $K_1(Z\times_kS)_{\mathbf Q}$ on the actual fixed
product. In particular every class satisfying
$\psi^2\eta=8\eta$ lifts, since $A(8)=168$.
Furthermore,
\begin{equation}
 \left\{x\in\ker\left(K_0(Z\times S)_{\mathbf Q}
                    \longrightarrow K_0(E\times S)_{\mathbf Q}\right)
                  :\psi^2x=8x\right\}=0.
                                                        \tag{PS.3}
\end{equation}
The latter is injectivity only on the indicated eigenspace.
Neither assertion supplies an extension across a nodal boundary
of a varying parameter curve.

\emph{1. Actual sections over the local base.}
Fix an ample line bundle $H$ on $S$ and pull it back to $B$.
For a vector bundle $V$, relative Serre vanishing applied to
the coherent sheaf $\mathfrak mV$ shows that, for sufficiently
large $n$, the closed-fibre restriction
\[
 \Gamma(B,VH^n)\longrightarrow\Gamma(S,(VH^n)|_S)
\]
is surjective \citep[Lemma 30.17.1]{frC3azStacksSerre}.
Such a twist on the smooth projective surface also generates
first jets. One direct uniform argument uses the finite flat
family of first infinitesimal neighbourhoods of the points of $S$:
the relative Serre vanishing theorem for its ideal makes the
evaluation onto every such neighbourhood surjective. Thus we
can choose a general closed-fibre section and lift it to the
actual product.

If $\operatorname{rk}V\geq3$, the incidence of sections vanishing
at a point of the surface has codimension at least three over
that point. Its image in the finite-dimensional section space
is closed, because $S$ is proper, and has dimension strictly
smaller than that space. We choose a section outside it.
The zero scheme of its lift to $B$ is proper over $R$ and misses
the closed fibre. A nonempty closed image in a local scheme
contains the closed point, so that zero scheme is empty.
The section gives a line subbundle with locally free quotient.
Iteration expresses every vector-bundle class as a sum of line
classes and one rank-two bundle class. We assert no splitting
of these extensions.

For a rank-two bundle, a general sufficiently twisted section
has finite transverse zero scheme on $S$. Here is the relevant
first-jet dimension estimate. At a specified point the two
values impose codimension two; failure of surjectivity of the
derivative is one further determinantal condition. The bad
incidence therefore has codimension at least three before the
point varies over the two-dimensional surface. Its proper
image is a proper closed subset of the section space.
Outside it, the zeros have independent relative derivatives.

Let $F$ be the zero scheme of the lifted section on $B$.
The relative Jacobian criterion makes $F\to\operatorname{Spec}R$
\'etale near every point of its closed fibre. Its non-\'etale
locus is closed in the proper $F$ and misses that fibre, hence
is empty. Properness and quasi-finiteness make $F$ finite
\'etale \citep[Lemma 37.44.1]{frNKStacksProperFinite}.
The same argument applies to two transverse high-twist line
sections. The individual divisors are relative effective
Cartier: the fibrewise criterion holds near the closed fibre,
and the proper-local argument covers the whole product
\citep[Lemma 31.19.9(3)]{frNKStacksRelativeCartier}.
All the finite zero schemes are regular immersions of
codimension two in the smooth relative surface $B/R$.
Their Koszul resolutions are therefore perfect even when $R$
is singular.

\emph{2. Finite supports annihilate the augmentation ideal.}
The product $B$ is connected; let $I$ be the rank-zero ideal of
$K_0(B)$. Let $G$ be the subgroup generated by the classes
$[\mathcal O_F]$ of the finite zero schemes above, including
the transverse intersections of pairs of line sections.
Every finite $R$-algebra is semilocal. A finite projective
module of constant rank over a semilocal ring is free:
choose bases at its finitely many maximal ideals, lift them
simultaneously by the Chinese remainder theorem, and use
Nakayama's lemma at every maximal ideal. The resulting map
from a free module is an isomorphism after every such
localization, and hence an isomorphism. For varying ranks
we apply the same argument on the open and closed rank
pieces. In particular all line bundles there are trivial.

Every class in $I$ restricts to zero in $K_0(F)$, on every
rank piece. If $j:F\hookrightarrow B$ is one of our perfect
immersions, the perfect projection formula gives
\begin{equation}
 x[\mathcal O_F]=j_*j^*x=0\quad(x\in I),
 \qquad IG=0.
                                                        \tag{PS.4}
\end{equation}
Tensoring a generator of $G$ by a line bundle leaves it
unchanged. The use of perfect immersions is essential;
no arbitrary coherent sheaf is being identified with a
perfect class. As $B$ has an ample line bundle over its
affine base, vector bundles compute its perfect $K_0$
\citep[Corollary 3.9 and Proposition 3.10]
{D_ThomasonTrobaugh1990}.

\emph{3. Exact degree-two generation.}
For any line bundles $M,N$ on $B$ we first prove
\begin{equation}
                         ([M]-1)([N]-1)\in G.
                                                        \tag{PS.5}
\end{equation}
Choose high twists $MH^a,H^a,NH^b,H^b$ and sections whose
closed-fibre divisors are smooth and transverse for each of
the four pairs used below. These are finitely many nonempty
open first-jet conditions, so they can be satisfied
simultaneously over the infinite field $k$.
Write their relative Cartier divisors as
$D_{M,+},D_{M,-},D_{N,+},D_{N,-}$.
The Cartier sequences give the exact identity
\begin{equation}
 [M]-1=[M][\mathcal O_{D_{M,+}}]
                         -[\mathcal O_{D_{M,-}}].
                                                        \tag{PS.6}
\end{equation}
Indeed its right side is
$[M](1-[M]^{-1}[H]^{-a})-(1-[H]^{-a})$.
Multiply (PS.6) by the corresponding formula for $N$.
The four products are the perfect Koszul classes of their
finite transverse intersections, with line twists.
Those twists are trivial on the semilocal intersections.
This proves (PS.5), including a square of a line difference
by using two independently chosen divisor representations.

For a rank-two bundle $V$, put
\[
 c(V)=\gamma^2([V]-2)=[\det V]-[V]+1,
 \qquad d=[\det V]-1.
\]
Choose $h=[H^n]$ so that $W=V\otimes H^n$ has one of the
transverse sections above. Its Koszul resolution gives
$[\mathcal O_F]=1-[W^*]+[\det W^*]$.
Multiplication by $\det W$ turns the right side into $c(W)$,
since $W^*\otimes\det W\simeq W$; it leaves the supported
class unchanged. Therefore $c(W)=[\mathcal O_F]$.
Direct expansion, retaining every term, gives
\begin{equation}
\begin{split}
 c(W)&=h\,c(V)+h(h-1)d+(h-1)^2,\\
 c(V)&=h^{-1}[\mathcal O_F]-(h-1)d-h^{-1}(h-1)^2\in G.
\end{split}                                               \tag{PS.7}
\end{equation}
The membership follows from (PS.5) and the line-twist
invariance of $G$.

Our bundle reduction and $[V]-2=d-c(V)$ now show that $I$
is generated by line differences and elements of $G$.
Thus $I^2\subseteq G$ and $I^3=0$, using (PS.4).
Conversely each defining finite support is a rank-two
Koszul class as above, so $G\subseteq F_\gamma^2K_0(B)$.
We still check the full gamma bound, rather than infer it
from augmentation nilpotence. For line and rank-two
generators respectively,
\[
 \gamma_z([M]-1)=1+([M]-1)z,\qquad
 \gamma_z([V]-2)=1+([V]-2)z+c(V)z^2.
\]
Gamma series of arbitrary virtual rank-zero classes are
products and inverses of these expressions. Every term
of total degree at least three contains either three
factors from $I$ or one from $G$ and one from $I$.
It therefore vanishes. The same argument applies to
products of gamma operations of total degree at least
three. We have proved both $F_\gamma^3K_0(B)=0$ and
$F_\gamma^2K_0(B)=G$, integrally.

\emph{4. Adams polynomial and the actual surface lift.}
The Adams operation $\psi^2$ acts by $1$ on the rank quotient,
by $2$ on $I/G$, and by $4$ on $G$. The last assertion can be
checked directly on a rank-two Koszul generator:
\[
 \psi^2c(V)=c(V)([\det V]+[V]+1)=4c(V),
\]
since the correction to the last bracket is in $I$ and
$IG=0$. The product of the three quotient-annihilating
operators gives $A(\psi^2)=0$ on all of $K_0(B)$.
No splitting or semisimplicity is assumed. Taking $R=k$
gives the same assertion for $K_0(S)$.

In characteristic $p$, FF.1--FF.2 gives the exact
fixed-product sequence
\begin{equation}
\begin{split}
 K_1(Z\times S)_{\mathbf Q}\oplus K_1(S)_{\mathbf Q}
 &\longrightarrow K_1(E\times S)_{\mathbf Q}\\
 &\longrightarrow K_0(B)_{\mathbf Q}\\
 &\longrightarrow K_0(Z\times S)_{\mathbf Q}
                         \oplus K_0(S)_{\mathbf Q}.
\end{split}                                               \tag{PS.8}
\end{equation}
Its inputs are the abstract-blow-up pro-exact comparison
\citep[Theorem A and (1.6)]{rootKerzStrunkTamme2018},
characteristic-$p$ nilpotent invariance after inverting $p$
\citep[Theorem A]{rootGeisserHesselholt2011}, and finite-cover
descent \citep[Theorem 8.1]{D_ThomasonTrobaugh1990}, at the
exact scopes proved in FF. Naturality makes the boundary
of $A(\psi^2)\eta$ zero by (PS.1). The $K_1(S)$ contribution
is already the restriction of its actual pullback to
$Z\times S$, proving (PS.2). If $\eta$ has eigenvalue $8$,
division by $168$ gives its full lift.

Finally, for $x$ in the kernel in (PS.3), the pair $(x,0)$
is in the image of $K_0(B)_{\mathbf Q}$ by exactness.
The polynomial therefore kills $x$. If $\psi^2x=8x$,
this says $168x=0$, giving (PS.3). Every map is on the
original fixed product. In particular a family with a
literal weight-three endpoint can first be transformed
by $A(\psi^2)$; its endpoint is multiplied by $168$ even
if the family itself is not a pure eigenspace class.
\qed

\emph{Scope and the remaining boundary obstruction.}
This is an auxiliary theorem. Spreading a varying
proper curve over $k(u)$ into a proper parameter surface
can leave a nodal boundary. Its supported localization
class need not be a sum of classes from the smooth
boundary components. Support-union exactness has
negative-$K$ terms at their intersections; the relevant
shift can retain $K_{-1}(E)_{\mathbf Q}$ of weight one.
Neither this theorem nor normalization of the boundary
removes those terms. We have not proved their vanishing
for the actual HK/VE construction, nor a replacement
which yields the required symbolic proper family.
The remaining UK.7 quotient and all main-row statuses
are consequently unchanged.
